\documentclass[10pt]{article}
\usepackage[top=0.75in,bottom=0.75in,left=0.75in,right=0.75in]{geometry}
\usepackage[T1]{fontenc}
\usepackage{enumitem}
\usepackage{titlesec}
\usepackage{array}
\usepackage{booktabs}

\title{AB-900: Copilot and Agent Administration Fundamentals}
\author{}
\date{}

\setcounter{secnumdepth}{2}
\setcounter{tocdepth}{2}
\renewcommand{\thesection}{\arabic{section}}
\renewcommand{\thesubsection}{\thesection.\arabic{subsection}}
\titleformat{\section}[display]{\normalfont\Huge\bfseries\centering}{\thesection}{0.5em}{\vspace*{\fill}}[\vspace*{\fill}]
\titlespacing*{\section}{0pt}{0pt}{0pt}
\titleformat{\subsection}[display]{\normalfont\LARGE\bfseries\centering}{\thesubsection}{0.5em}{\vspace*{1.25cm}}[\vspace{0.6em}\titlerule]
\titlespacing*{\subsection}{0pt}{0pt}{1.5em}

\begin{document}
\begin{titlepage}
\thispagestyle{empty}
\vspace*{\fill}
\begin{center}
{\LARGE\bfseries AB-900: Copilot and Agent Administration Fundamentals\par}
\end{center}
\vspace*{\fill}
\end{titlepage}
\tableofcontents
\newpage

\setcounter{section}{0}
\section{Microsoft 365 Security Foundations}
\setcounter{subsection}{0}
\clearpage
\subsection{Analyze the Zero Trust Security Model}

\begin{center}
\begin{tabular}{>{\raggedright\arraybackslash}p{0.28\textwidth} >{\raggedright\arraybackslash}p{0.6\textwidth}}
\toprule
\textbf{Abbreviation} & \textbf{Full Form} \\
\midrule
PIM & Privileged Identity Management \\
BYOD & Bring-your-own-device \\
DLP & Data Loss Prevention \\
NIST & National Institute of Standards and Technology \\
ISO & International Organization for Standardization \\
\bottomrule
\end{tabular}
\end{center}

\begin{enumerate}[leftmargin=*]

\item In Microsoft 365, the Zero Trust model is best described as:
\begin{enumerate}[label=\alph*., leftmargin=1.5em]
\item A single product
\item A comprehensive strategy that spans identity, endpoints, data, applications, infrastructure, and networks
\item A replacement for all compliance programs
\item A feature available only in Microsoft Defender for Endpoint
\end{enumerate}
\textbf{Answer:} b. A comprehensive strategy that spans identity, endpoints, data, applications, infrastructure, and networks.

\textbf{Explanation:} The documentation explicitly says the model isn't a single product but a comprehensive strategy across six domains.

\item Which native Microsoft tools are specifically named as part of implementing Zero Trust in Microsoft 365? \textbf{(Select all that apply)}
\begin{enumerate}[label=\alph*., leftmargin=1.5em]
\item Microsoft Entra ID
\item Microsoft Defender
\item Intune
\item Purview
\item Windows Server Update Services
\end{enumerate}
\textbf{Answer:} a, b, c, d.

\textbf{Explanation:} The documentation says implementing Zero Trust requires a layered approach combining native tools like Microsoft Entra ID, Microsoft Defender, Intune, and Purview.

\item According to the documentation, why are organizations having a Zero Trust conversation? \textbf{(Select all that apply)}
\begin{enumerate}[label=\alph*., leftmargin=1.5em]
\item IT security is complex
\item Trusted network security strategy accepted lower security within the network
\item Assets increasingly leave the network
\item Attackers shift to identity attacks
\item Passwords are no longer used anywhere
\end{enumerate}
\textbf{Answer:} a, b, c, d.

\textbf{Explanation:} The documentation lists these four reasons in the Zero Trust conversation diagram.

\item The three core principles of the Zero Trust security model are Verify explicitly, Least privilege access, and Assume breach. \textbf{(True/False)}

\textbf{Answer:} True.

\textbf{Explanation:} These are the three principles introduced in the documentation under Core Principles of Zero Trust.

\item Which principle is described as the cornerstone of the Zero Trust model?
\begin{enumerate}[label=\alph*., leftmargin=1.5em]
\item Least privilege access
\item Assume breach
\item Verify explicitly
\item Endpoint compliance
\end{enumerate}
\textbf{Answer:} c. Verify explicitly.

\textbf{Explanation:} The documentation says the Verify explicitly principle is the cornerstone of the Zero Trust model.

\item Verify explicitly requires that every access request is authenticated and authorized using all available contextual signals. \textbf{(True/False)}

\textbf{Answer:} True.

\textbf{Explanation:} The documentation explains that every request from a user, device, or application must be evaluated using contextual signals.

\item Which contextual signals are named in the documentation? \textbf{(Select all that apply)}
\begin{enumerate}[label=\alph*., leftmargin=1.5em]
\item Identity attributes
\item Device health
\item Location
\item User behavior
\item Risk level
\end{enumerate}
\textbf{Answer:} a, b, c, d, e.

\textbf{Explanation:} All five are explicitly listed as contextual signals under Verify explicitly.

\item Microsoft 365 enforces Verify explicitly through which combination of mechanisms? \textbf{(Select all that apply)}
\begin{enumerate}[label=\alph*., leftmargin=1.5em]
\item Conditional Access policies
\item Real-time risk evaluation
\item Multifactor authentication
\item Static credentials only
\end{enumerate}
\textbf{Answer:} a, b, c.

\textbf{Explanation:} The documentation says these mechanisms work together to make access decisions dynamic and based on current context rather than static credentials.

\item Conditional Access in Microsoft Entra ID enables administrators to define:
\begin{enumerate}[label=\alph*., leftmargin=1.5em]
\item Granular access rules based on multiple signals
\item Only password reset workflows
\item Only firewall rules
\item Data labeling templates
\end{enumerate}
\textbf{Answer:} a. Granular access rules based on multiple signals.

\textbf{Explanation:} Conditional Access evaluates user identity, device compliance, location, application sensitivity, and real-time risk before granting access.

\item Which scenario is specifically mentioned in the documentation as a Conditional Access example?
\begin{enumerate}[label=\alph*., leftmargin=1.5em]
\item Require MFA for users accessing SharePoint Online from outside the corporate network
\item Allow all personal devices to download sensitive documents
\item Disable compliance requirements for contractors
\item Remove access controls for trusted locations permanently
\end{enumerate}
\textbf{Answer:} a. Require MFA for users accessing SharePoint Online from outside the corporate network.

\textbf{Explanation:} This is one of the exact examples given in the Conditional Access section.

\item Conditional Access also supports which session controls? \textbf{(Select all that apply)}
\begin{enumerate}[label=\alph*., leftmargin=1.5em]
\item Limiting access to read-only mode
\item Preventing downloads
\item Enforcing data watermarking in every app
\item Blocking all sign-ins from corporate devices
\end{enumerate}
\textbf{Answer:} a, b.

\textbf{Explanation:} The documentation specifically lists read-only mode and preventing downloads as session controls.

\item Policies can be combined using logical operators to create complex rules that reflect organizational requirements. \textbf{(True/False)}

\textbf{Answer:} True.

\textbf{Explanation:} The documentation explains that logical operators can combine multiple conditions such as device compliance, trusted location, and low risk score.

\item Secure Score integrates with Conditional Access to do what?
\begin{enumerate}[label=\alph*., leftmargin=1.5em]
\item Highlight gaps and suggest improvements
\item Replace Identity Protection
\item Create DLP labels automatically
\item Disable policy simulations
\end{enumerate}
\textbf{Answer:} a. Highlight gaps and suggest improvements.

\textbf{Explanation:} The documentation says Secure Score helps organizations implement Conditional Access effectively by surfacing gaps and recommendations.

\item Administrators can simulate policy impact before deployment to avoid unintended disruptions. \textbf{(True/False)}

\textbf{Answer:} True.

\textbf{Explanation:} The documentation highlights simulation as a way to keep verification rigorous and user-friendly.

\item Microsoft Entra Identity Protection evaluates sign-in risk using:
\begin{enumerate}[label=\alph*., leftmargin=1.5em]
\item Machine learning and global threat intelligence
\item Hardware inventory only
\item Local group policy only
\item Bandwidth monitoring only
\end{enumerate}
\textbf{Answer:} a. Machine learning and global threat intelligence.

\textbf{Explanation:} The documentation explicitly says sign-in risk is evaluated using machine learning and global threat intelligence.

\item Which anomalies are specifically named under risk-based authentication? \textbf{(Select all that apply)}
\begin{enumerate}[label=\alph*., leftmargin=1.5em]
\item Impossible travel
\item Unfamiliar sign-in properties
\item Known compromised credentials
\item Missing document labels
\end{enumerate}
\textbf{Answer:} a, b, c.

\textbf{Explanation:} These are the specific anomalies listed under Microsoft Entra Identity Protection.

\item Based on risk level, access can be:
\begin{enumerate}[label=\alph*., leftmargin=1.5em]
\item Blocked
\item Challenged with MFA
\item Allowed with restrictions
\item All of the above
\end{enumerate}
\textbf{Answer:} d. All of the above.

\textbf{Explanation:} The documentation says access can be blocked, challenged with MFA, or allowed with restrictions depending on the risk level.

\item MFA adds a critical layer of security by requiring users to verify their identity using a second factor. \textbf{(True/False)}

\textbf{Answer:} True.

\textbf{Explanation:} The documentation describes a second factor such as a mobile app notification, biometric scan, or hardware token.

\item Which MFA methods or passwordless options are specifically mentioned? \textbf{(Select all that apply)}
\begin{enumerate}[label=\alph*., leftmargin=1.5em]
\item Mobile app notification
\item Biometric scan
\item Hardware token
\item Windows Hello
\item FIDO2 keys
\end{enumerate}
\textbf{Answer:} a, b, c, d, e.

\textbf{Explanation:} All five are named in the documentation as supported verification methods or passwordless options.

\item Risk-based MFA can be configured to trigger only when necessary, minimizing user friction. \textbf{(True/False)}

\textbf{Answer:} True.

\textbf{Explanation:} The documentation explains this design as a balance between security and usability.

\item Scenario: A user signs in from a known device and location. Based on the documentation, what might happen?
\begin{enumerate}[label=\alph*., leftmargin=1.5em]
\item The user might not be prompted for MFA
\item The user must always be blocked
\item The user must always be prompted for hardware token verification
\item The sign-in is automatically marked high risk
\end{enumerate}
\textbf{Answer:} a. The user might not be prompted for MFA.

\textbf{Explanation:} The documentation says a sign-in from a known device and location might not trigger MFA, while a new device in a high-risk location can require it.

\item Scenario: A user accesses sensitive data from a corporate laptop in a trusted location. What outcome best matches the documentation?
\begin{enumerate}[label=\alph*., leftmargin=1.5em]
\item Seamless access might be allowed
\item Access must always be denied
\item Global administrator access must be granted
\item Device health is ignored
\end{enumerate}
\textbf{Answer:} a. Seamless access might be allowed.

\textbf{Explanation:} The documentation uses this as a lower-risk scenario compared to access from a personal device in a foreign location.

\item What does the Least privilege access principle dictate?
\begin{enumerate}[label=\alph*., leftmargin=1.5em]
\item Users, applications, and services should be granted only the minimum level of access necessary to perform their tasks
\item All administrators should have permanent privileged access
\item Every compliant device should receive full access
\item Access should be based only on job seniority
\end{enumerate}
\textbf{Answer:} a. Users, applications, and services should be granted only the minimum level of access necessary to perform their tasks.

\textbf{Explanation:} The documentation states this principle reduces the attack surface and contains damage if an account is compromised.

\item Implementing least privilege requires a thorough understanding of what? \textbf{(Select all that apply)}
\begin{enumerate}[label=\alph*., leftmargin=1.5em]
\item User roles
\item Workflows
\item Resource sensitivity
\item Wallpaper settings
\end{enumerate}
\textbf{Answer:} a, b, c.

\textbf{Explanation:} The documentation says organizations must map out users, workflows, and resource sensitivity to avoid overprovisioning.

\item Which capabilities support least privilege in Microsoft 365? \textbf{(Select all that apply)}
\begin{enumerate}[label=\alph*., leftmargin=1.5em]
\item RBAC
\item PIM
\item Defender for Cloud Apps only
\item Azure Firewall only
\end{enumerate}
\textbf{Answer:} a, b.

\textbf{Explanation:} RBAC provides predefined permission roles, and PIM adds just-in-time access to sensitive functions.

\item Which example is used to explain RBAC?
\begin{enumerate}[label=\alph*., leftmargin=1.5em]
\item A helpdesk technician granted the Password Administrator role
\item A contractor granted Global Administrator permanently
\item A guest user granted mailbox access for all employees
\item A mobile phone granted domain admin rights
\end{enumerate}
\textbf{Answer:} a. A helpdesk technician granted the Password Administrator role.

\textbf{Explanation:} The documentation uses Password Administrator as an example of a narrowly defined role that avoids privilege creep.

\item Match each control or product with its described purpose.

\begin{center}
\begin{tabular}{>{\raggedright\arraybackslash}p{0.35\textwidth} >{\raggedright\arraybackslash}p{0.55\textwidth}}
\toprule
\textbf{Column A} & \textbf{Column B} \\
\midrule
1. RBAC & a. Aggregates logs and correlates events across domains \\
2. PIM & b. Provides just-in-time access to sensitive functions \\
3. Microsoft Sentinel & c. Defines roles with specific permissions \\
4. Microsoft Purview Information Protection & d. Classifies, labels, and encrypts sensitive data \\
\bottomrule
\end{tabular}
\end{center}

\textbf{Answer:} 1-c, 2-b, 3-a, 4-d.

\textbf{Explanation:} Each product or control is mapped directly to the role described in the documentation.

\item In PIM, a global administrator can request temporary elevation for a specific task, such as configuring a new Conditional Access policy. \textbf{(True/False)}

\textbf{Answer:} True.

\textbf{Explanation:} The documentation gives this exact example when explaining PIM and just-in-time elevation.

\item Once a privileged task is complete, the elevated privileges are automatically revoked. \textbf{(True/False)}

\textbf{Answer:} True.

\textbf{Explanation:} Automatic revocation is one of the reasons PIM reduces the risk of persistent high-level access.

\item The Assume breach principle reflects a shift in mindset from prevention to containment. \textbf{(True/False)}

\textbf{Answer:} True.

\textbf{Explanation:} The documentation says Assume breach operates under the assumption that attackers are already inside or might eventually get in.

\item What does the Assume breach principle emphasize? \textbf{(Select all that apply)}
\begin{enumerate}[label=\alph*., leftmargin=1.5em]
\item Segmentation
\item Continuous monitoring
\item Rapid response
\item Permanent trust for internal systems
\end{enumerate}
\textbf{Answer:} a, b, c.

\textbf{Explanation:} These three actions are explicitly named as the focus areas for minimizing breach impact.

\item Which tools are specifically named under the Assume breach principle? \textbf{(Select all that apply)}
\begin{enumerate}[label=\alph*., leftmargin=1.5em]
\item Microsoft Defender for Endpoint
\item Microsoft Sentinel
\item Microsoft Defender for Identity
\item Azure Policy
\end{enumerate}
\textbf{Answer:} a, b, c.

\textbf{Explanation:} These are the three security tools listed as providing visibility across endpoints, identities, and network traffic.

\item Microsoft Defender for Endpoint uses behavioral analytics to detect which anomalies? \textbf{(Select all that apply)}
\begin{enumerate}[label=\alph*., leftmargin=1.5em]
\item Lateral movement
\item Privilege escalation
\item Ransomware activity
\item SharePoint site creation
\end{enumerate}
\textbf{Answer:} a, b, c.

\textbf{Explanation:} These are the exact behaviors named in the Defender for Endpoint section.

\item If a device exhibits suspicious behavior, Microsoft Defender for Endpoint can automatically isolate it from the network. \textbf{(True/False)}

\textbf{Answer:} True.

\textbf{Explanation:} The documentation says isolation helps prevent further spread of suspicious activity.

\item What is Microsoft Sentinel described as doing?
\begin{enumerate}[label=\alph*., leftmargin=1.5em]
\item Aggregating logs and alerts from across Microsoft 365 and non-Microsoft sources
\item Managing only local firewalls
\item Enforcing app protection policies on mobile phones
\item Replacing Conditional Access
\end{enumerate}
\textbf{Answer:} a. Aggregating logs and alerts from across Microsoft 365 and non-Microsoft sources.

\textbf{Explanation:} The documentation says Sentinel uses machine learning to correlate events and identify threats spanning multiple domains.

\item Which scenario is used to explain Microsoft Sentinel correlation?
\begin{enumerate}[label=\alph*., leftmargin=1.5em]
\item A phishing email detected by Defender for Office 365 linked to a compromised account in Microsoft Entra ID and a malicious file in SharePoint
\item A printer jam linked to a failed password reset
\item A Teams status message linked to a VM resize request
\item A DLP label linked to a device wallpaper change
\end{enumerate}
\textbf{Answer:} a. A phishing email detected by Defender for Office 365 linked to a compromised account in Microsoft Entra ID and a malicious file in SharePoint.

\textbf{Explanation:} This is the exact cross-domain example used to describe Sentinel’s correlation capabilities.

\item Which techniques can Microsoft Defender for Identity detect? \textbf{(Select all that apply)}
\begin{enumerate}[label=\alph*., leftmargin=1.5em]
\item Pass-the-Hash
\item Golden Ticket attacks
\item Domain enumeration
\item Browser theme changes
\end{enumerate}
\textbf{Answer:} a, b, c.

\textbf{Explanation:} These identity attack techniques are explicitly listed in the documentation.

\item Assume breach also supports proactive threat hunting and incident response. \textbf{(True/False)}

\textbf{Answer:} True.

\textbf{Explanation:} The documentation says security teams can use telemetry data to trace attack paths, identify compromised assets, and contain threats.

\item What are the six interdependent Zero Trust pillars in Microsoft 365?

\textbf{Answer:} Identity, Endpoints, Data, Applications, Infrastructure, and Network.

\textbf{Explanation:} The documentation names these six pillars as the domains where Zero Trust principles must be applied.

\item Which pillar is described as the cornerstone of Zero Trust?
\begin{enumerate}[label=\alph*., leftmargin=1.5em]
\item Identity
\item Applications
\item Data
\item Infrastructure
\end{enumerate}
\textbf{Answer:} a. Identity.

\textbf{Explanation:} The documentation explicitly states that identity is the cornerstone of Zero Trust.

\item Microsoft Entra ID provides which core identity capabilities? \textbf{(Select all that apply)}
\begin{enumerate}[label=\alph*., leftmargin=1.5em]
\item Centralized identity management
\item Authentication
\item Authorization
\item Remote imaging
\end{enumerate}
\textbf{Answer:} a, b, c.

\textbf{Explanation:} These three are named in the Identity section as core capabilities of Microsoft Entra ID.

\item Which passwordless sign-in examples are specifically listed? \textbf{(Select all that apply)}
\begin{enumerate}[label=\alph*., leftmargin=1.5em]
\item Windows Hello for Business
\item FIDO2 security keys
\item Microsoft Authenticator
\item FTP shared passwords
\end{enumerate}
\textbf{Answer:} a, b, c.

\textbf{Explanation:} These are the exact passwordless examples used in the Identity section.

\item Microsoft Entra ID supports identity federation, allowing users to sign in with credentials from external providers like Google or Facebook while still enforcing corporate policies. \textbf{(True/False)}

\textbf{Answer:} True.

\textbf{Explanation:} This sentence closely follows the wording used in the documentation.

\item In the context of the Zero Trust model, endpoints refers to what?

\textbf{Answer:} Any device that connects to enterprise resources, including desktops, laptops, mobile phones, tablets, and even IoT devices.

\textbf{Explanation:} The documentation defines endpoints broadly and notes that they are often the first target in cyberattacks.

\item Which tools are specifically highlighted for endpoint management and protection? \textbf{(Select all that apply)}
\begin{enumerate}[label=\alph*., leftmargin=1.5em]
\item Microsoft Intune
\item Microsoft Defender for Endpoint
\item Microsoft Purview Information Protection
\item Azure Policy
\end{enumerate}
\textbf{Answer:} a, b.

\textbf{Explanation:} The Endpoints section specifically names Intune and Defender for Endpoint.

\item Microsoft Intune enables administrators to define and enforce compliance policies that can require what? \textbf{(Select all that apply)}
\begin{enumerate}[label=\alph*., leftmargin=1.5em]
\item Encryption
\item Antivirus protection
\item Minimum OS versions
\item Other security configurations
\end{enumerate}
\textbf{Answer:} a, b, c, d.

\textbf{Explanation:} All four are listed as examples of compliance policy requirements.

\item Scenario: A user attempts to access Microsoft Teams from a jailbroken iPhone. Based on the documentation, what is the likely result?
\begin{enumerate}[label=\alph*., leftmargin=1.5em]
\item The user might be blocked due to noncompliance
\item The user is automatically granted access
\item Device compliance checks are skipped
\item The device is trusted because it is mobile
\end{enumerate}
\textbf{Answer:} a. The user might be blocked due to noncompliance.

\textbf{Explanation:} This is the specific example used in the Intune section.

\item What does Endpoint Analytics provide?
\begin{enumerate}[label=\alph*., leftmargin=1.5em]
\item Insights into device performance, user behavior, and configuration risks
\item Email encryption only
\item Passwordless sign-in only
\item A replacement for Conditional Access
\end{enumerate}
\textbf{Answer:} a. Insights into device performance, user behavior, and configuration risks.

\textbf{Explanation:} Endpoint Analytics is described as helping identify crashes, outdated software, and insecure settings.

\item Which Microsoft tool provides visibility into sanctioned and unsanctioned applications, also known as shadow IT?
\begin{enumerate}[label=\alph*., leftmargin=1.5em]
\item Microsoft Defender for Cloud Apps
\item Microsoft Defender for Cloud
\item Microsoft Intune
\item Microsoft Purview Information Protection
\end{enumerate}
\textbf{Answer:} a. Microsoft Defender for Cloud Apps.

\textbf{Explanation:} The Applications section explicitly identifies Defender for Cloud Apps as the tool for shadow IT visibility and control.

\item Microsoft Entra ID’s App Proxy enables secure remote access to on-premises applications without exposing them to the internet. \textbf{(True/False)}

\textbf{Answer:} True.

\textbf{Explanation:} The documentation states that App Proxy uses reverse proxy architecture and integrates with Conditional Access.

\item Microsoft Purview Information Protection allows organizations to do what?
\begin{enumerate}[label=\alph*., leftmargin=1.5em]
\item Classify, label, and encrypt sensitive data
\item Restrict outbound traffic to known safe destinations
\item Assign helpdesk roles
\item Create VM size restrictions
\end{enumerate}
\textbf{Answer:} a. Classify, label, and encrypt sensitive data.

\textbf{Explanation:} This is the primary function described in the Data pillar.

\item A document containing credit card numbers can be automatically labeled as ``Confidential'' and encrypted. \textbf{(True/False)}

\textbf{Answer:} True.

\textbf{Explanation:} The documentation uses this exact example to explain automatic labeling and encryption.

\item DLP policies prevent unauthorized sharing of sensitive information. \textbf{(True/False)}

\textbf{Answer:} True.

\textbf{Explanation:} The documentation defines DLP as a control to block, warn, or allow with justification when sensitive information is shared.

\item DLP integrates with which Microsoft 365 services? \textbf{(Select all that apply)}
\begin{enumerate}[label=\alph*., leftmargin=1.5em]
\item Exchange
\item SharePoint
\item OneDrive
\item Teams
\end{enumerate}
\textbf{Answer:} a, b, c, d.

\textbf{Explanation:} These are the exact services listed in the documentation.

\item Microsoft Defender for Cloud supports multicloud environments, including which platforms?
\begin{enumerate}[label=\alph*., leftmargin=1.5em]
\item Amazon Web Services and Google Cloud
\item Only Microsoft Azure
\item Only on-premises datacenters
\item Only Microsoft 365 workloads
\end{enumerate}
\textbf{Answer:} a. Amazon Web Services and Google Cloud.

\textbf{Explanation:} The Infrastructure section explicitly says Defender for Cloud supports multicloud environments including AWS and Google Cloud.

\item Which compliance standards are specifically mentioned in the Infrastructure section? \textbf{(Select all that apply)}
\begin{enumerate}[label=\alph*., leftmargin=1.5em]
\item ISO 27001
\item National Institute of Standards and Technology (NIST)
\item PCI DSS
\item CIS Controls
\end{enumerate}
\textbf{Answer:} a, b.

\textbf{Explanation:} The documentation specifically mentions ISO 27001 and NIST.

\item What is Azure Policy described as?

\textbf{Answer:} A powerful governance tool that enables organizations to create, assign, and manage policies that enforce rules and effects across cloud resources.

\textbf{Explanation:} The Infrastructure section introduces Azure Policy as the tool used to enforce consistent governance and prevent misconfigurations.

\item Which examples are given for Azure Policy enforcement? \textbf{(Select all that apply)}
\begin{enumerate}[label=\alph*., leftmargin=1.5em]
\item Restrict VM sizes
\item Enforce tagging
\item Require encryption
\item Automatically isolate suspicious endpoints
\end{enumerate}
\textbf{Answer:} a, b, c.

\textbf{Explanation:} These are the exact governance examples listed under Azure Policy.

\item Network segmentation and access control are essential in Zero Trust. \textbf{(True/False)}

\textbf{Answer:} True.

\textbf{Explanation:} The Network section begins by stating this directly.

\item Which technologies are explicitly mentioned in the Network pillar? \textbf{(Select all that apply)}
\begin{enumerate}[label=\alph*., leftmargin=1.5em]
\item Azure Firewall
\item VPN Gateway
\item Microsoft Defender for Identity
\item Microsoft Authenticator
\end{enumerate}
\textbf{Answer:} a, b, c.

\textbf{Explanation:} These are the named technologies used to monitor and control network traffic and detect lateral movement.

\item Scenario: An organization wants to restrict outbound traffic to known safe destinations to prevent data exfiltration. Which control best matches the documentation?
\begin{enumerate}[label=\alph*., leftmargin=1.5em]
\item Azure Firewall
\item Endpoint Analytics
\item Password Administrator role
\item App Proxy
\end{enumerate}
\textbf{Answer:} a. Azure Firewall.

\textbf{Explanation:} The documentation specifically says Azure Firewall can restrict outbound traffic to known safe destinations.

\item Defender for Identity monitors network traffic for signs of lateral movement, such as credential theft or unauthorized access. \textbf{(True/False)}

\textbf{Answer:} True.

\textbf{Explanation:} This statement closely follows the wording used in the Network section.

\item Put the three core Zero Trust principles in the order presented in the documentation.

\textbf{Answer:} Verify explicitly, Least privilege access, Assume breach.

\textbf{Explanation:} The documentation introduces the principles in this exact order.

\item Scenario: A phishing email leads to credential theft. Which response best matches the layered defense strategy described in the documentation?
\begin{enumerate}[label=\alph*., leftmargin=1.5em]
\item Defender for Identity detects lateral movement, while Microsoft Sentinel correlates the event with other indicators of compromise
\item Intune automatically changes all document labels to Confidential
\item Azure Policy revokes all sign-in sessions for every user
\item Endpoint Analytics creates a new Conditional Access policy automatically
\end{enumerate}
\textbf{Answer:} a. Defender for Identity detects lateral movement, while Microsoft Sentinel correlates the event with other indicators of compromise.

\textbf{Explanation:} The documentation uses this example to show how Assume breach supports proactive threat hunting and incident response through layered defenses.

\end{enumerate}

\end{document}
